\documentclass[11pt]{article}
%% \usepackage{epsfig}
%% \usepackage{color}  
%% \usepackage{oldlfont}
%\usepackage[scanall]{psfrag}
%% \usepackage{graphics}
%% \usepackage{amssymb}
%% \usepackage{amsmath}
\usepackage{hyperref}
\addtolength{\topmargin}{-1cm}
\addtolength{\textheight}{2cm}
\addtolength{\textwidth}{4cm}
\addtolength{\oddsidemargin}{-2cm}
\begin{document}


\subsection*{Introduction}
The Beach-fx modeling system requires application of a
one-dimensional coastal morphology change model that is suitably
accurate for predicting beach profile change under storm attack.  The
original framework made use of the SBEACH model, an empirical
representation of morphodynamics introduced in Larson and Kraus (1989).
The option to use the CSHORE model was outlined in Johnson and
Sanderson (2020) and the scripting code is provided \\
\href{https://github.com/erdc/BeachFX_CSHORE_link}{https://github.com/erdc/BeachFX\_CSHORE\_link}\\
Both models require comparable bottom position initial conditions with
wave and water level seaward boundary data.  In addition to profile
response, SBEACH and CSHORE also provide hydrodynamic and wave 
estimates required by the Beach-fx model.  While the models are
generally interchangeable, they are contrasting in foundational
principles and associated limiting assumptions.  A short review of the
operational model differences is provided herein to assist in rational
model choice.



\subsection*{Model Formulation} 

Both SBEACH and CSHORE operate under the assumption of longshore
uniformity, and wave and current estimates derive from the numerical
integration of the depth-averaged energy, momentum, and continuity
equations.

The predictions of sediment transport and associated morphology change
are fundamentally different, however.  The SBEACH model pre-determines
an empirically-based equilibrium profile on the basis of sediment, wave and
water level conditions.  The rate of bottom position evolution,
approaching the prescribed steady-state, is based on the degree of
departure from equilibrium and the hydrodynamic forcing.  With a basis
in prescribed steady profiles, The SBEACH model lacks foundational rigor
but is inherently robust and stable.


The CSHORE model, conversely, posits process-based sediment transport
algorithms for a nearshore breaking wave environment. The model
provides distinct estimates for bedload, suspended load, and
wave-related sediment transport.  Generally speaking, when applied in
an appropriate manner, the CSHORE model provides more accurate results
[e.g. Johnson etal 2012].  The underlying assumptions for the CSHORE
model are, however, more restrictive, and predictions resulting from
conditions in violation will likely be in gross error.  A partial
accounting of model assumptions to be regarded is:
\begin{itemize}
\item
  Sediment grain size varies in the range $.1 mm <d_{50}<1 mm$,
\item
  Minimal bed position change at the seaward boundary, equivalent to
  having a seaward boundary sufficiently deep to avoid breaking,
\item
  Dune overtopping is limited to intermittent wave overtopping.  This
  precludes having conditions where the seaward boundary water level
  is higher than the maximum profile position.
\item 
 Slopes are limited to the sediment angle of repose.  This precludes
 the use of vertical walls, which are not currently supported within
 CSHORE.
\end{itemize}

\subsection*{Model Operation} 

At present, the SBEACH workflow is designed to execute on a Windows
platform in serial.  It is common practice to manually subdivide a
long series of Monte Carlo simulations and perform the morphology
computations on separate individual computers and manually recombine
the results at completion.  In contrast, the CSHORE/BEACHFX
implementation is scripted for rudimentary parallel execution, where
an arbitrary number of discreet processors or cores can be used
simultaneously on a desktop computer or HPC resources.  While the HPC
platform can pose challenges, the parallel system affords considerable
efficiency.


\subsection*{References} 

\noindent Johnson, B.D. Sanderson, D.R. (2020) {\it On the Use of CSHORE for
  Beach-fx}, US Army Corps of Engineers Technical Note CHETN-II-59\\[1mm]

\noindent Johnson, B.D., Kobayashi, N., and Gravens, M.B. (2012) {\it
  Cross-Shore Numerical Model CSHORE for Waves, Currents, Sediment
  Transport and Beach Profile Evolution}, US Army Corps of Engineers
Technical Report TR-12-22.\\[1mm]

\noindent Larson, M., Kraus, N.C. (1989) {\it SBEACH : numerical model for
  simulating storm-induced beach change. Report 1, Empirical
  foundation and model development} US Army Corps of Engineers
Technical Report; CERC-89-9 Report 1



\end{document}
